% preamble.tex —— 《迹公式与函子性：一个纲领的开端》中译注本 全局导言
% 编译：xelatex ×2（有索引时 xelatex → makeindex → xelatex ×2）
\usepackage[a4paper,margin=2.6cm,footskip=1.2cm]{geometry}
% 页锚放在页边：book 类默认 \marginparwidth=116pt，比 2.6cm 的页边距还宽，
% marginnote 会被排到纸外；这里把它压到页边距以内。
\setlength{\marginparwidth}{1.9cm}
\setlength{\marginparsep}{2pt}
% 中西文混排 + 长公式常造成 overfull，放宽断行
\setlength{\emergencystretch}{3em}
\tolerance=2500
\usepackage{ctex}
\usepackage{amsmath,amssymb,amsthm,mathtools}
\usepackage{mathrsfs,bbm}
\usepackage{stmaryrd}
\usepackage{tikz}
\usepackage{tikz-cd}
\usetikzlibrary{arrows.meta,decorations.pathmorphing,calc}
\usepackage{xcolor}
\usepackage{enumitem}
\usepackage{booktabs,array,multirow}
\usepackage{graphicx}
\usepackage{marginnote}
\usepackage[unicode,bookmarksnumbered,bookmarksopen=true]{hyperref}
\usepackage{imakeidx}
\makeindex[intoc]

% ---------- 译者体例开关：\showyizhufalse 可一键隐藏全部译注 ----------
\newif\ifshowyizhu
\showyizhutrue

% ---------- 颜色 ----------
\definecolor{yizhuGray}{gray}{0.30}
\definecolor{srcBlue}{RGB}{20,70,140}

% ---------- 译注：\yizhu{内容}{来源} ----------
% 体例：以脚注形式排，楷体，来源用 [来源：…] 标注；无外部依据写 [译者按（无依据）]
\newcommand{\yizhu}[2]{%
  \ifshowyizhu
    \footnote{\textcolor{yizhuGray}{\kaishu #1}\quad{\scriptsize\textcolor{srcBlue}{[来源：#2]}}}%
  \fi}
\newcommand{\yizhuw}[1]{%
  \ifshowyizhu
    \footnote{\textcolor{yizhuGray}{\kaishu #1}\quad{\scriptsize\textcolor{srcBlue}{[译者按（无外部依据）]}}}%
  \fi}

% ---------- 原文页码锚点：\srcpage{7} ----------
\newif\ifsrcmark
\srcmarktrue
\newwrite\srcpageout
\immediate\openout\srcpageout=\jobname.pagemarks
\newcommand{\srcpage}[1]{%
  \ifsrcmark\immediate\write\srcpageout{#1}\fi
  \phantomsection
  \marginnote{\raggedleft\scriptsize\textcolor{yizhuGray}{p.#1}}}
% 无 margin 的精简版（用于窄页/表格内）：\pgref{N}
\newcommand{\pgref}[1]{\textsuperscript{\scriptsize\textcolor{yizhuGray}{[p.#1]}}}
% 页界落在句子中间时用这个：行内小号上标，保证不会被排版吞掉
\newcommand{\srcpageinl}[1]{%
  \ifsrcmark\immediate\write\srcpageout{#1}\fi
  \phantomsection
  \textsuperscript{\tiny\textcolor{yizhuGray}{〔p.#1〕}}}

% ---------- 编号：单篇论文，节号从 1 起 ----------
\renewcommand{\sectionmark}[1]{\markright{\thesection\quad #1}}
\renewcommand{\thesection}{\arabic{section}}
\renewcommand{\thesubsection}{\thesection.\arabic{subsection}}
\renewcommand{\thesubsubsection}{\thesection.\arabic{subsection}.\arabic{subsubsection}}

% ---------- 定理环境：编号按「节.序」，与原文 Lemme 2.6、Proposition 3.29 对齐 ----------
\theoremstyle{plain}
\newtheorem{theorem}{定理}[section]
\newtheorem{proposition}[theorem]{命题}
\newtheorem{lemma}[theorem]{引理}
\newtheorem{corollary}[theorem]{推论}
\theoremstyle{definition}
\newtheorem{definition}[theorem]{定义}
\newtheorem{example}[theorem]{例}
\newtheorem{conjecture}[theorem]{猜想}
\theoremstyle{remark}
\newtheorem{remark}[theorem]{注}

% ---------- 常用记号 ----------
\newcommand{\FF}{\mathbb{F}}
\newcommand{\QQ}{\mathbb{Q}}
\newcommand{\ZZ}{\mathbb{Z}}
\newcommand{\RR}{\mathbb{R}}
\newcommand{\CC}{\mathbb{C}}
\newcommand{\PP}{\mathbb{P}}
\renewcommand{\AA}{\mathbb{A}}% 覆盖 Å
\newcommand{\GG}{\mathbb{G}}
\newcommand{\oo}{\mathcal{O}}
\newcommand{\cB}{\mathcal{B}}
\newcommand{\cD}{\mathcal{D}}
\newcommand{\cE}{\mathcal{E}}
\newcommand{\cF}{\mathcal{F}}
\newcommand{\cG}{\mathcal{G}}
\newcommand{\cH}{\mathcal{H}}
\newcommand{\cL}{\mathcal{L}}
\newcommand{\cM}{\mathcal{M}}
\newcommand{\cO}{\mathcal{O}}
\newcommand{\cT}{\mathcal{T}}
\newcommand{\cX}{\mathcal{X}}
\newcommand{\cZ}{\mathcal{Z}}
\newcommand{\Bun}{\mathrm{Bun}}
\newcommand{\Hck}{\mathrm{Hecke}}
\newcommand{\Hecke}{\mathrm{Hecke}}
\newcommand{\shF}{\mathcal F}


\newcommand{\Orb}{\mathrm{Orb}}
\newcommand{\Rep}{\mathrm{Rep}}
\newcommand{\GL}{\mathrm{GL}}
\newcommand{\SL}{\mathrm{SL}}
\newcommand{\Gal}{\mathrm{Gal}}
\newcommand{\Frob}{\mathrm{Frob}}
\newcommand{\Spec}{\operatorname{Spec}}
\newcommand{\Aut}{\operatorname{Aut}}
\newcommand{\Hom}{\operatorname{Hom}}
\newcommand{\End}{\operatorname{End}}
\newcommand{\Sym}{\operatorname{Sym}}
\newcommand{\Tr}{\operatorname{Tr}}
\newcommand{\rk}{\operatorname{rk}}
\newcommand{\Lie}{\operatorname{Lie}}
\newcommand{\id}{\mathrm{id}}
\newcommand{\im}{\operatorname{im}}
\newcommand{\coker}{\operatorname{coker}}
\newcommand{\divi}{\operatorname{div}}
\newcommand{\res}{\operatorname{res}}
\newcommand{\nrd}{N}
\newcommand{\Ql}{\overline{\QQ}_\ell}
\newcommand{\LG}{{}^{L}G}
\newcommand{\LH}{{}^{L}H}
\newcommand{\lamG}{{}^{\lambda}G}
\newcommand{\lamH}{{}^{\lambda}H}
\newcommand{\lamGphi}{{}^{\lambda}G_{\varphi}}
\newcommand{\lamHpi}{{}^{\lambda}H_{\pi}}
\newcommand{\wt}[1]{\widetilde{#1}}
\newcommand{\ol}[1]{\overline{#1}}
\newcommand{\ul}[1]{\underline{#1}}
\newcommand{\p}{\mathfrak{p}}
\newcommand{\g}{\mathfrak{g}}
\newcommand{\sG}{\boldsymbol{G}}
\DeclareMathOperator{\tam}{tam}
\DeclareMathOperator{\meas}{meas}
\DeclareMathOperator{\cnd}{cnd}

% 交叉引用
\newcommand{\thmref}[1]{定理~\ref{#1}}
\newcommand{\propref}[1]{命题~\ref{#1}}
\newcommand{\lemref}[1]{引理~\ref{#1}}
\newcommand{\defref}[1]{定义~\ref{#1}}
\newcommand{\remref}[1]{注~\ref{#1}}
